\section{Translation attacks on frozen detectors and text watermarks}
\label{sec:translation_experiments}

\paragraph{Question and protocol.}
Does generating outside English and translating back into English remove
detector evidence or a watermark? The generator is Qwen3.5-9B instruct, in
three versions: unwatermarked, the native TextSeal~2 sampler, and the local
DistSeal-inspired latent text student. The last is a merged model distilled for
5,000 updates, not an official released DistSeal text checkpoint. All versions
receive an English news prefix and an instruction to continue in Simplified
Chinese, French or Welsh. Index-Translate-9B and TranslateGemma-12B independently
translate each saved continuation into English, without adding a watermark.
This is foreign-language generation followed by translation, not a round trip
starting from a saved English answer.

The fixed cohort is separate from the CNN/DailyMail warmup and validation IDs
described in Section~\ref{sec:adaptation}. The Chinese comparison uses
the first 2,000 test prefixes per version, and the French/Welsh pilot uses the
same first 200. Source generation uses temperature 1, top-$p$ 0.95 and a
256-new-token cap; translation is greedy with a 1,024-token cap. Token budgets
are tokenizer-specific, so they do not imply equal information across
languages. Saved output prefixes and checkpointed inference allow exact
resumption without regenerating completed texts.

\paragraph{Separate detector and watermark outcomes.}
The primary AI detector is the unchanged equal-weight mean of the seven
Qwen-warmed member scores, with threshold 0.5. No translated texts are used
for fitting or recalibration. TextSeal~2 detection uses its native
$p\leq 0.001$ rule. The text student's extractor threshold, approximately
$-0.615143$, is frozen from human and unwatermarked English validation
responses. A source watermark is scored before translation, and the same
owner test is applied afterward. Conditional watermark loss is the fraction
of \emph{initially detectable source watermarks} no longer detected after
translation; it is different from the raw translated detection rate.

\paragraph{Chinese-to-English results.}
Direct-English ensemble detection is 99.80\% for the clean generator,
99.95\% for TextSeal~2 and 98.35\% for the text student. The corresponding
English watermark detection rates are 99.60\% and 58.65\%. On the
Chinese-requested sources, watermark detection is 1,989/2,000 (99.45\%) for
TextSeal~2 and only 308/2,000 (15.40\%) for the student. The student's weak
source-language baseline must therefore be separated from translation-induced
loss. Table~\ref{tab:translation_chinese} reports final English outcomes.
Human false positives on the matched 2,000 English controls are 0.50\% for
the ensemble, 0.05\% for TextSeal~2 and 0.60\% for the student.

\begin{table}[ht]\centering\small
\begin{adjustbox}{max width=\linewidth}
\begin{tabular}{lllll}\toprule
Version & Translator & AI detection & WM detected & WM loss$^*$\\\midrule
Clean & Index & 99.75\% & N/A & N/A\\
Clean & TranslateGemma & 99.75\% & N/A & N/A\\
TextSeal 2 & Index & 99.70\% & 7.05\% & 92.91\%\\
TextSeal 2 & TranslateGemma & 99.75\% & 5.50\% & 94.47\%\\
Text student & Index & 99.10\% & 7.60\% & 55.52\%\\
Text student & TranslateGemma & 99.35\% & 5.85\% & 66.56\%\\
\bottomrule\end{tabular}\end{adjustbox}
\caption{Chinese-requested generation followed by translation into English, 2,000 outputs per condition. AI detection uses the frozen ensemble at 0.5; WM detected is the translated watermark detection rate. $^*$WM loss conditions on a detected source watermark: 1,989 TextSeal 2 sources and 308 text-student sources. Loss is not a content-preserving removal rate. Unwatermarked controls have no WM test.}\label{tab:translation_chinese}\end{table}


\paragraph{French and Welsh pilot.}
Table~\ref{tab:translation_french_welsh} keeps source language, model version
and translator separate. At the fixed ensemble threshold, the first 200 human
controls have 3/200 false positives (1.50\%); this differs from the larger
Chinese cohort's empirical rate. The clean-model French outputs are all
detected after either translator, whereas Welsh-to-English through Index
shows greater evasion. The TextSeal~2 French-source watermark is detectable in
200/200 sources, falling to 5/200 after Index and 4/200 after TranslateGemma.
The Welsh TextSeal~2 sources have 199/200 detectable watermarks.

The released DivEye component requires at least 15 words. One 12-word
Welsh-to-English translation stopped the initial job. The resumed protocol
preserves every output, records attempted and scorable counts, and treats
unsupported short texts as unscorable for the full ensemble. It does not
substitute a smaller ensemble or manufacture a score. A valid-evasion count
requires at least 20 words and uses \emph{all attempted outputs} as its
denominator. Watermark detection retains the full output set independently
of detector coverage.

\begin{table}[ht]\centering\small
\begin{adjustbox}{max width=\linewidth}
\begin{tabular}{lllllll}\toprule
Version & Source & Source WM & AI Index & AI Gemma & WM Index & WM Gemma\\\midrule
Clean & French & N/A & 100.00\% & 100.00\% & N/A & N/A\\
Clean & Welsh & N/A & 89.50\% & 97.00\% & N/A & N/A\\
TextSeal 2 & French & 100.00\% & 98.00\% & 98.50\% & 2.50\% & 2.00\%\\
TextSeal 2 & Welsh & 99.50\% & 90.50\% & 98.49\%$^{\dagger}$ & 9.00\% & 7.00\%\\
Text student & French & 4.50\% & 97.00\% & 97.50\% & 3.50\% & 2.50\%\\
Text student & Welsh & 21.50\% & 82.00\% & 95.00\% & 11.50\% & 10.00\%\\
\bottomrule\end{tabular}\end{adjustbox}
\caption{French/Welsh pilot: 200 attempted outputs per condition; Gemma means TranslateGemma. Source WM is measured before translation. AI columns use the same frozen seven-member ensemble and 0.5 threshold. Watermark columns retain all 200 outputs. Texts below the released DivEye minimum of 15 words are unscorable for the full ensemble and are not successful valid evasions. $\dagger$TextSeal 2/Welsh/TranslateGemma: 199/200 scorable.}\label{tab:translation_french_welsh}\end{table}


\paragraph{Interpretation and limits.}
The observations separate two kinds of evidence: the frozen ensemble largely
retains discrimination in the Chinese and French conditions, while translation
removes most detectable TextSeal~2 marks; Welsh is more sensitive to the chosen
translator. They do not establish universal translation robustness or
content-preserving watermark removal. Language instructions are imperfectly
followed: 226/2,000 clean and 146/2,000 TextSeal~2 Chinese-requested sources
contain no Han characters, with manually inspected examples in English.
These are raw requested-language results, not a fully language-validated
Chinese subset. Source-language compliance and translation adequacy remain
unassessed for the complete matrix, including Welsh. Student thresholds were
calibrated on English rather than language-matched human controls.

The direct-English quality check uses 2,000 samples and custom Mistral
final-token MAUVE features with three clustering seeds. Values are low across
versions: 0.0205--0.0240 for the clean model, 0.0210--0.0232 for TextSeal~2,
and 0.0377--0.0420 for the student. These custom-feature scores are not directly
comparable to default GPT-2 MAUVE. This proxy alone does not establish preserved quality, and it is
not a translated-output adequacy test. The smaller pilots also have limited
resolution and uncertainty. Portable metrics, source hashes, coverage and
frozen protocol identities are retained in
\path{tex/provenance/translation_transfer_results.json}. No translation-robust
detector or watermark defense is trained in these experiments.
